\documentclass[11pt]{article}

\usepackage[T1]{fontenc}
\usepackage{lmodern}
\usepackage{amsmath,amssymb,amsthm}
\usepackage{booktabs}
\usepackage{enumitem}
\usepackage{microtype}
\usepackage[hidelinks]{hyperref}

\newtheorem{theorem}{Theorem}
\newtheorem{proposition}[theorem]{Proposition}
\newtheorem{lemma}[theorem]{Lemma}
\theoremstyle{definition}
\newtheorem{definition}[theorem]{Definition}

\newcommand{\NAE}{\operatorname{NAE}}
\newcommand{\Sat}{\operatorname{Sat}}
\newcommand{\Fthree}{\mathbb F_{3}}

\title{An Explicit 98-Clause Linear Monotone NAE-4 Equality Gadget}
\author{Xiaojian Zeng\\
  \small PBC School of Finance, Tsinghua University, Beijing 100190, China\\
  \small \texttt{xiaojianzeng@foxmail.com}}
\date{}

\begin{document}

\maketitle

\begin{abstract}
We give an explicit bounded-size equality gadget using only monotone
four-variable not-all-equal clauses, with any two clauses meeting in at most
one variable.  The gadget is obtained by deleting one line from an explicitly
specified 99-line non-two-colourable subhypergraph of
$\operatorname{PG}(3,3)$.  Its terminal projection is exactly
$\{0000,1111\}$, as certified by a dependency-free exhaustive verifier.
Replacing every indexed occurrence in a monotone NAE-4 instance by a fresh
copy and attaching one gadget therefore preserves the full projected solution
set and produces a linear instance.  An input with $n$ variables and $m$
clauses is transformed into one with $n+156m$ variables and $393m$ clauses in
linear time; the local bound may be taken as $B=98$.  This supplies the
four-uniform linearization step invoked, but not exhibited, in the reduction
of Bhardwaj (arXiv:2609.03508v1).  No claim is made about minimality of the
gadget or about unrelated downstream steps of that reduction.
\end{abstract}

\section{Context and statement}

An indexed monotone NAE-4 instance is a pair
\[
   F=(V;(C_j)_{j=1}^{m}), \qquad C_j\in \binom{V}{4}.
\]
The indices matter, so identical members of $\binom{V}{4}$ may occur at
different indices.  A map $s:V\to\{0,1\}$ satisfies a clause when its four
values are not all equal.  The instance is \emph{linear} when
$|C_i\cap C_j|\leq 1$ for distinct indices $i,j$.

Bhardwaj's reduction from restricted NAE-4-SAT to recognition of Max-Cut SDP
exactness invokes a standard vertex-splitting linearization, tied together by
intermediate NAE equality clauses \cite[Sec.~3.1, Lemma~1]{Bhardwaj2026}.
The subsequent $K_6$ construction uses linearity to ensure that its clause
cliques are edge-disjoint \cite[Sec.~3.2]{Bhardwaj2026}.  The cited proof does
not specify a four-uniform equality gadget or its terminal projection.

The closest explicit construction we found is due to Darmann and D\"ocker.
Their 12-clause linear monotone NAE-3 gadget forces four terminals to have a
common value, and it is used occurrence-by-occurrence to linearize a
3-uniform formula \cite[Lemma~3 and Theorem~2]{DarmannDocker2019}.  Because
all of its clauses have arity three, it does not by itself provide the
four-uniform replacement needed here.  Our contribution is limited to the
following explicit four-uniform construction and its composition proof; it
is not a priority claim about all possible NAE-4 gadgets.

\begin{theorem}\label{thm:main}
For every finite indexed monotone NAE-4 instance
$F=(V;(C_j)_{j=1}^{m})$, there is a monotone NAE-4 instance $F'$ such that:
\begin{enumerate}[label=\textup{(\roman*)}]
  \item every two distinct clauses of $F'$ intersect in at most one variable;
  \item
  \[
     \{s'|_V:s'\in\Sat(F')\}=\Sat(F);
  \]
  in particular, every satisfying assignment of $F$ extends to one of $F'$
  and every satisfying assignment of $F'$ restricts to one of $F$;
  \item $F'$ has $n+156m$ variables and $393m$ clauses, where $n=|V|$;
  and
  \item $F'$ is generated in $O(n+m)$ time from fixed incidence data.
\end{enumerate}
Each replaced occurrence uses 39 fresh variables and 98 new clauses, so the
same absolute bound $B=98$ applies to both local quantities.
\end{theorem}

\section{The projective equality gadget}

We first specify the finite incidence structure without relying on a stored
list of point coordinates or clauses.  Represent the 40 points of
$\operatorname{PG}(3,3)$ by the vectors in $\Fthree^4$ whose first nonzero
coordinate is $1$, ordered lexicographically, and label them $0,\ldots,39$.
For every linearly independent pair of representatives, normalize the four
one-dimensional subspaces in its two-dimensional span.  Deduplicate the
resulting four-element point sets, sort each set by point label, and order the
130 sets lexicographically.  Call the resulting lines
$\ell_0,\ldots,\ell_{129}$; this deterministic convention gives
\[
  \ell_0=\{0,1,2,3\}=:L.
\]

Let $K$ be the hypergraph on these 40 points whose line-index set is
\begingroup\small
\begin{align*}
I=\{&0,2,3,4,5,6,7,8,9,11,12,13,17,18,19,20,21,22,24,26,27,28,30,31,\\
&33,34,35,37,38,39,40,41,43,44,46,47,49,50,51,52,54,56,57,60,61,\\
&63,64,65,66,67,68,70,71,73,74,75,76,77,80,81,82,83,84,85,88,91,\\
&92,93,94,95,96,97,98,99,100,101,102,103,105,106,108,109,110,111,\\
&112,113,114,115,116,117,118,119,120,121,122,123,127,128,129\}.
\end{align*}
\endgroup
Define the 98-edge hypergraph
\[
   G:=K\setminus\{L\}.
\]
Every edge of $G$ has four distinct points.  Moreover, two distinct
projective lines are two-dimensional vector subspaces whose intersection
cannot contain two distinct one-dimensional subspaces unless the lines are
equal.  Thus $G$ is a linear 4-uniform hypergraph.

\begin{proposition}[Exact terminal projection]\label{prop:projection}
The NAE-satisfying colourings of $G$, projected onto the ordered terminal set
$(0,1,2,3)$, are exactly
\[
   \{(0,0,0,0),(1,1,1,1)\}.
\]
\end{proposition}

\begin{proof}
Two exact finite facts are used:
\begin{enumerate}[label=\textup{(\alph*)}]
  \item $K$ has no NAE two-colouring;
  \item $G$ has a colouring whose four terminal values are zero.
\end{enumerate}
They are certified exhaustively as described below.  Complementing the
colouring in (b) gives one with terminal pattern $1111$.  Conversely, if a
colouring of $G$ had a nonmonochromatic terminal pattern, it would also
satisfy the deleted NAE edge $L$, giving a colouring of $K$ and contradicting
(a).  These are therefore the only two terminal patterns.
\end{proof}

Identify terminal 0 with a variable $x$, terminal 1 with a variable $y$, and
regard terminals 2 and 3 and the remaining 36 points as auxiliary variables.
Proposition~\ref{prop:projection} then says precisely that an assignment on
$\{x,y\}$ extends through $G$ if and only if $x=y$.  This is stronger than
equisatisfiability: it determines the entire two-terminal projection.

\section{Exact finite certificate}\label{sec:certificate}

The accompanying program \texttt{gadget\_certificate.py} constructs the
points and lines from arithmetic in $\Fthree$, selects $I$, and runs a complete
two-colour search.  Its only propagation rule is forced by a four-edge: if
three assigned vertices of an edge have the same colour and its fourth vertex
is unassigned, the fourth receives the opposite colour.  A completed
monochromatic edge is a conflict.  When neither condition applies, the
program selects an unassigned vertex and recursively tries both colours.

This recursion is complete by induction on the number of unassigned
vertices.  Propagation discards no satisfying colouring, because the forced
value is the only value that can satisfy the relevant edge.  At a branch,
the two recursive calls partition all extensions of the current partial
assignment.  Hence a return value of UNSAT after both calls is an exhaustive
statement, not a floating-point or heuristic conclusion.

For $K$, the run visits 15,047 search nodes, reports 7,524 conflicts and
105,948 propagations, and returns UNSAT.  For $G$ with its terminals fixed to
zero it returns the following 40-bit witness in point order:
\[
\texttt{0000000111001001011001100111001010101011}.
\]
Direct substitution checks every one of the 98 edges.  The program also fixes
each terminal pattern and exhausts the remaining variables, producing the
full local table
\begin{center}
\small
\begin{tabular}{@{}cccccccc@{}}
\toprule
0000 & 0001 & 0010 & 0011 & 0100 & 0101 & 0110 & 0111\\
SAT  & UNSAT& UNSAT& UNSAT& UNSAT& UNSAT& UNSAT& UNSAT\\
\midrule
1000 & 1001 & 1010 & 1011 & 1100 & 1101 & 1110 & 1111\\
UNSAT& UNSAT& UNSAT& UNSAT& UNSAT& UNSAT& UNSAT& SAT\\
\bottomrule
\end{tabular}
\end{center}

In certified finite searches it is useful to separate the correctness of an
encoding, termination of a solver, and checking of the resulting proof object;
Zeng, Liu, Ratnavelu, and Huat make this distinction explicitly in a SAT-based
geometric setting \cite[Sec.~1]{ZengEtAl2026}.  Here the fixed instance is
small enough that the supplied dependency-free verifier performs the complete
branching itself.  A separate verification pass also regenerated all 40
points and 130 lines, encoded each NAE edge by the two CNF clauses forbidding
all-zero and all-one values, and reproduced the UNSAT and 16-row projection
results with an independent SAT backend.  The theorem, however, depends only
on the inspectable incidence convention, the exhaustive verifier, and the
completeness argument above.

\section{Occurrence replacement and composition}

Write
\[
   C_j=\{x_{j,1},x_{j,2},x_{j,3},x_{j,4}\},
\]
where the four variables in a clause are distinct.  For every indexed
occurrence $(j,r)$, introduce a fresh variable $y_{j,r}$ and replace $C_j$ by
the copied clause
\[
   D_j=\{y_{j,1},y_{j,2},y_{j,3},y_{j,4}\}.
\]
Attach a fresh copy $G_{j,r}$ of $G$ by identifying its terminal 0 with the
original variable $x_{j,r}$ and its terminal 1 with $y_{j,r}$.  All other
vertices of all gadget copies are fresh.  This replaces every occurrence,
which in particular replaces every offending occurrence.

\begin{lemma}\label{lem:linear}
The resulting instance $F'$ is linear.
\end{lemma}

\begin{proof}
There are five exhaustive types of clause pairs.  Two copied input clauses
are disjoint because all occurrence variables are fresh.  Two clauses inside
one gadget meet in at most one variable by projective linearity.  Gadget
copies attached to different original variables are disjoint.  Two gadget
copies attached to the same original variable share only that variable.
Finally, a copied input clause and a gadget clause can share only the unique
occurrence variable belonging to that gadget.  Every pair therefore has
intersection size at most one.
\end{proof}

\begin{lemma}\label{lem:projection}
The restriction map from $\Sat(F')$ to assignments on $V$ has image exactly
$\Sat(F)$.
\end{lemma}

\begin{proof}
Let $s\in\Sat(F)$.  Set $y_{j,r}=s(x_{j,r})$.  When this common value is zero,
use the displayed zero-terminal witness in $G_{j,r}$; when it is one, use its
bitwise complement.  These choices are independent between gadget copies.
Each $D_j$ then receives the same four values as $C_j$, so all copied clauses
and all gadget clauses are satisfied.  Thus $s$ extends to $F'$.

Conversely, let $s'\in\Sat(F')$.  Proposition~\ref{prop:projection} applied to
each copy $G_{j,r}$ gives
$s'(y_{j,r})=s'(x_{j,r})$.  Since $D_j$ is satisfied, substituting these four
equalities shows that $C_j$ is satisfied by $s'|_V$.  This holds for every
indexed clause, so $s'|_V\in\Sat(F)$.
\end{proof}

For additional finite checks of the composition boundary, the verifier forms
pairs of indexed input clauses sharing respectively two, three, and four
variables.  The transformed instances have respectively 318, 317, and 316
variables, each has 786 clauses, and in every case the maximum intersection
of two distinct output clauses is one.  Exhaustion of the 64, 32, and 16
input truth-table rows also confirms that the canonical gadget extension
exists exactly on the original satisfying rows.  These tests are not needed
to replace Lemmas~\ref{lem:linear} and \ref{lem:projection}; they exercise all
possible offending overlap sizes in the implementation.

\begin{proof}[Proof of Theorem~\ref{thm:main}]
Lemmas~\ref{lem:linear} and \ref{lem:projection} prove the incidence and
assignment assertions.  Each gadget copy has 40 vertices, but terminal 0 is
identified with an existing original variable.  It therefore introduces 39
fresh variables, including the occurrence variable, and 98 clauses.  There
are $4m$ indexed occurrences, giving
\[
   |V(F')|=n+39(4m)=n+156m
\]
and
\[
   |C(F')|=m+98(4m)=393m.
\]
The line-index set and all normalized incidences are fixed data independent
of the input.  Creating a constant-size gadget for each occurrence and
copying the $m$ input clauses takes $O(n+m)$ time.  Since $39\leq98$, one may
take the requested common local bound to be $B=98$.
\end{proof}

\section{Scope and limitations}

The result establishes the bounded four-uniform replacement, linearity,
extension, restriction, and complexity assertions for every finite indexed
monotone NAE-4 instance with four distinct variables per clause.  It does not
use signed literals, constants, repeated variables within a clause, or
lower-arity constraints.  The transformation preserves the projected
solution set, not just the existence of a solution.

We do not assert that 98 clauses or 39 fresh variables are minimal.  The
finite core was found computationally, and the paper makes no unrestricted
priority claim beyond the stated comparison with the inspected sources.  The
construction fills the linearization step used by Bhardwaj's unweighted
$K_6$ reduction under the monotone four-uniform interpretation above.  It
does not audit the later $K_6$ calculation, the weighted reduction, membership
claims, or any other conclusion in that paper.

\paragraph{Computational and AI assistance.}
A language model assisted with the search organization, literature triage,
and preparation of the exposition.  The fixed finite claims were checked by
the dependency-free exhaustive program described in
Section~\ref{sec:certificate}; a separate verification pass reconstructed the
incidence system and checked the Boolean assertions with an independent SAT
backend.  The author is responsible for the statement and proof presented
here.

\paragraph{Code availability.}
The supplementary files \texttt{gadget\_certificate.py} and
\texttt{gadget\_certificate.out} reproduce the projective incidence data,
core non-two-colourability, terminal projection, witness check, linearity
check, and overlap-composition tests using exact integer and Boolean
arithmetic.

\bibliographystyle{alpha}
\bibliography{references}

\end{document}
